\documentclass{article}
\usepackage{graphicx}
\usepackage{amsmath}
\usepackage{amssymb}
\usepackage{amsfonts}
\usepackage{hyperref}
\usepackage{exercises}
\usepackage{indentfirst}
\usepackage{circledsteps}
\usepackage{amsthm}
\usepackage{csquotes}
\usepackage{yhmath}
\usepackage{comment}
\usepackage{placeins}
\usepackage{geometry}
\usepackage{hyperref}
\usepackage{tikz}
\newcommand*{\oldneg}{\mathord{\sim}}
\newcommand{\Z}{\mathbb{Z}}
\newcommand{\N}{\mathbb{N}}
\let\oldemptyset\emptyset

\theoremstyle{definition}
\newtheorem{definition}{Definition}
\newtheorem{theorem}{Theorem}[section]
\newtheorem{corollary}{Corollary}[theorem]
\newtheorem{lemma}[theorem]{Lemma}

\setlength{\parindent}{.5in}
\geometry{margin=0.5in}
\setlength{\parskip}{12pt}
\title{Chapter 9 Notes}
\begin{document}
\date{}
\maketitle

\section*{Section 1}

\subsection*{Notes}
Here is why $A\cup B=A\cup (B-A)$.

On the one hand, suppose that $x\in A\cup B$. Then either $x\in A$ or $x\in B$. In the first case, obviously 
$x\in A\cup (B-A)$. In the second case, if $x\in A$, then we are done (by the first case). Otherwise $x\in B-A$, meaning 
$x\in A\cup (B-A)$. On the other hand, suppose that $x\in A\cup (B-A)$. Then obviously $x\in A\cup B$.

\section*{Section 2}

\subsection*{Notes}
Note that Lemma 2.2 assumes that $h$ is defined on all subsets of $A$.

In proof 2 of the Bernstein-Schröder theorem, here is why $A-[A-g''(B-f''(X))]=g''(B-f''(X))$. 
Obviously, if $x\in A-[A-g''(B-f''(X))]$, then $x\in g''(B-f''(X))$. On the other hand, if 
$x\in g''(B-f''(x))$, note that $x\in A$, as $g$ is a function from $B$ into $A$. Thus 
$x\notin A-g''(B-f''(X))$, meaning $x\in A-[A-g''(B-f''(X))]$.

\subsection*{Exercise 2.1}
(a) Consider any fixed point $A$ of $h$. Then $h(A)=A$, meaning $A\subseteq h(A)$. Hence $A\in\Sigma$, and thus 
$A\subseteq B$. Thus $B$ is the largest fixed point of $h$.

(b) Let $\Sigma$ be the set of all subsets $X$ of $A$ such that $X\subseteq h(X)$ and that $X$ is a subset of all fixed 
points of $h$. Let $B=\cup\Sigma$. We aim to show that $B$ is a fixed point of $h$. Obviously, $B$ is a subset of all fixed 
points of $h$ (thus $B$ is the smallest fixed point of $h$).
	
	On the one hand, we show that $B\subseteq h(B)$ (this part is identical to the proof of the first part of the original lemma). 
So suppose that $b\in B$. Then $b\in X$ for some $X\in\Sigma$.
Because $X\in\Sigma$, then $X\subseteq h(X)$, hence $b\in h(X)$. But $X\subseteq B$, and because $h$ is monotone, 
it follows that $h(X)\subseteq h(B)$. Thus $x\in h(B)$. Hence $B\subseteq h(B)$.
	
	On the other hand, we show that $h(B)\subseteq B$. Well, we have shown that $B\subseteq h(B)$, meaning 
$h(B)\subseteq h(h(B))$ (because $h$ is monotone). Also, consider any fixed point $A$ of $h$. Then $B\subseteq A$. 
Hence $h(B)\subseteq h(A)$. But $h(A)=A$, meaning $h(B)\subseteq A$. Thus $h(B)$ is a subset of every fixed point of 
$h$, and $h(B)\subseteq h(h(B))$, meaning $h(B)\in\Sigma$, and thus $h(B)\subseteq B$.
	Hence we have shown that $h(B)=B$, meaning $B$ is a fixed point of $h$ (and the smallest one).

\section*{Section 3}

\begin{lemma}
For any classes $A$ and $B$, $A\prec B$ and $B\prec A$ do not both hold.
\end{lemma}
\begin{proof}
This is immediate from the Bernstein-Schröder theorem. For suppose $A\prec B$ and $B\prec A$ both held. 
Then $A\preceq B$ and $B\preceq A$ both held, meaning, by the Bernstein-Schröder theorem, $A\cong B$. But this contradicts 
the fact that $A\cong B$ was assumed to not hold (as $A\prec B$ was assumed to hold).
\end{proof}

\begin{lemma}
For any class $A$, no subclass of $A$ is greater than $A$.
\end{lemma}
\begin{proof}
Suppose a subclass $B$ of $A$ was greater than $A$. Then $A\prec B$ holds. Also 
$B\prec A$ holds (this is because $B\preceq A$ holds as $B\subseteq A$, but
$B\cong A$ does not hold by assumption). But this contradicts Lemma 0.1. Hence 
$B$ is not greater than $A$.
\end{proof}

\begin{lemma}
All subsets of any countable set are countable.
\end{lemma}
\begin{proof}
Recall that, by Exercise 7.1 of Chapter 3, every set of natural numbers is either finite or denumerable. Now 
consider any countable set $A$. Either it is finite or denumerable.

(1) On the one hand, suppose that $A$ is finite. Then (by Exercise 6.1 of Chapter 3) all subsets of $A$ are also
finite.

(2) On the other hand, suppose that $A$ is denumerable. Then it can be enumerated in some 
sequence $a_0,a_1,...,a_n,...$. Now consider any subset $S$ of $A$. Then 
$S$ can be put into a 1-1 correspondence with the set $C$ of all $n$ such that 
$a_n\in S$. But $C$ is a set of natural numbers, meaning it is countable. Thus $S$ is also 
countable.
\end{proof}

\begin{lemma}
The set of all sets of natural numbers with $n$ as its greatest element is finite, in particular, it has 
$2^n$ elements.
\end{lemma}
\begin{proof}
By induction. Obviously the base case holds (as there is only one set of natural numbers with $0$ as its greatest element). 
Suppose the result holds for some $n$. 
Now, the set of all sets of 
natural numbers with $n+1$ as its greatest element is just all sets of the form $\cup\{\{n+1\},x\}$, where $x\in\mathcal{P}(n)$. But 
there are $2^{n+1}$ elemetns in $\mathcal{P}(n)$. Hence the set of all sets of natural numbers with $n+1$ as its greatest element 
has $2^{n+1}$ elements.
\end{proof}

\begin{lemma}
The union of a denumerable set of disjoint finite sets is denumerable.
\end{lemma}
\begin{proof}
Suppose we have a denumerable set $S$ of disjoint finite sets $F_0,F_1,...,F_n,...$. For each $F_i$, let 
$N(F_i)$ be the size of $F_i$. Then define $f_0$ to be a function which assigns the elements of 
$F_0$ to the first $N(F_i)$ numbers. Then, for each $n$, let $f_{n^+}$ be a function containing the ordered pairs of 
$f_n$, and if $f_n$ contains a finite number of elements, add to $f_{n^+}$ the ordered pairs which assigns the elements of 
$F_{n^+}$ to the first $N(F_{n^+})$ elements of $\omega-\textrm{card}(f_n)$. If $f_n$ is infinite, then 
take $f_{n^+}=f_n$. An obvious induction argument shows that each $f_n$ must be finite, meaning another easy 
induction argument shows that each $f_n$ has all elements of each $F_i$ for $i\le n$ in its domain. Then take the function 
$f$ to be the set of all ordered pairs $\langle x,f_n(x)\rangle$ where $x\in F_n$. An easy induction argument shows that 
$f$ is a 1-1 function between $\cup S$ and $\omega$.
\end{proof}

\subsection*{Exercise 3.1}

(1) Consider any denumerable set $A$. It can be placed in some sequence $a_0,a_1,...,a_n,...$. Then $A$ can be put into a 1-1 correspondence 
with $A-\{a_0\}$, a proper subset of $A$.

(2) Consider some finite set $D$ of denumerable sets. It can be placed in some sequence $a_0,...,a_n$. For each 
$a_i$, let $a'_i$ be the set of all elements $x$ of $a_i$ such $i$ is the least index of $D$ that contains $x$ (that is, 
if $x\in a_j$ and $j\ne i$, then $j>i$). Obviously 
$\cup S=\cup\{a'_0,...,a'_n\}$, and all elements $a'_i$ are disjoint with all other elements $a'_j$, where $i\ne j$. Each 
$a_i$ is denumerable, and each $a'_i\subseteq a_i$, meaning (by Lemma 0.3) that each $a'_i$ is countable. Thus each $a'_i$ is either 
finite or denumerable. Then the union of the set $F$ of all finite $a'_i$'s is finite. Then there are only a finite number of 
possible $a'_i$'s that are denumerable. Say the set $S$ of remaining $a'_i$'s can be enumerated by the sequence $b_0,...,b_m$, and, for each 
$b_i$, $b_i$ can be enumerated as $c^i_0,c^i_1,...,c^i_n,...$. Then, for each $k$, the set $S_k$ of all $c^i_k$ (for all $i\le m$) 
for every $k$ is finite. Also, 
$\cup S=\{S_0,S_1,...,S_n,...\}$, which is denumerable by Lemma 0.5. Thus $\cup D=(\cup F)\cup (\cup S)$, which is denumerable by 
part (5) of this exercise.

(3) Suppose that $A$ and $B$ are denumerable. Then $A$ can be placed in some sequence $a_0,a_1,...,a_n,...$ and $B$ can be placed in some 
sequence $b_0,b_1,...,b_n,...$. Then, for each $i$, the set $A_i$ of all ordered pairs $\langle a_i,b_j\rangle$ of $A\times B$ with $i\ge j$ is finite. 
Similarly, for each $i$, the set $B_i$ of all ordered pairs $\langle a_j,b_i\rangle$ of $A\times B$ with $i\ge j$ is finite. For all $i$ and $j$ where 
$i\ne j$, $A_i\cup B_i$ is disjoint from $A_j\cup B_j$. Then 
$A\times B=\cup\{A_0\cup B_0,A_1\cup B_1,...,A_n\cup B_n,...\}$, which is denumerable by Lemma 0.5.

(4) Suppose that $A$ is denumerable. Then it can be placed in some sequence $a_0,a_1,...,a_n,...$. Then the set of all finite subsets of $A$ with $i$ as its greatest index
is finite, in particular, it has $2^i$ elements (the proof of this is identical to the proof of Lemma 0.4). 
For each $i$, let $S(a_i)$ be the set containing all subsets of $A$ with $a_i$ as its greatest element. 
Then the set of all finite subsets of $A$ is just the set $\cup\{\{\emptyset\},S(a_0),S(a_1),...,S(a_n),...\}$ which is denumerable by Lemma 0.4. 

(5) Suppose that $A$ is denumerable and $F$ is finite. Say $A$ can be enumerated in the sequence $a_0,a_1,...,a_n,...$ and $F$ can be enumerated in the sequence $b_0,...,b_m$. 
Then, if $F-A$ has $k$ elements (it must be finite), we pair of the first $k$ numbers with the elements of $F-A$, and then the rest of the numbers can be placed in a 1-1 correspondence 
with $A$ itself.

\section*{Section 4}

\subsection*{Exercise 4.1}
Suppose that $A$ is infinite and there is a choice function $C$ for $A$. For any proper subset $x$ of $A$, 
let $g(x)=x\cup\{C(A-x)\}$, and let $g(A)=A$. By Theorem 8.1 of Chapter 3, 
there is a function $h$ from $\omega$ into $\mathcal{P}(A)$ such that $h(0)=g(0)$, and for any $n$, $h(n^+)=g(h(n))$. Then take $a_0$ to be the sole element of of $h(0)$ and for 
each $n$, take $a_{n^+}$ to be the sole element of $h(n^+)-h(n)$. Now, 
$a_0$ is the sole element of $h(0)=g(0)=0\cup\{C(A-0)\}=0\cup\{C(A)\}=\{C(A)\}$, that is, $a_0$ is the $C$-chosen element of $A$. Also note that 
$h(n^+)=g(h(n))=h(n)\cup\{C(A-h(n))\}$, that is, $h(n^+)$ is the set of all elements of $h(n)$, as well as the $C$-chosen element of $A-h(n)$. We aim to show that 
$a_{n^+}$ is the $C$-chosen element of $A-\{a_0,...,a_n\}$. First we show that $h(n)=\{a_0,...,a_n\}$. We have already proven the base case. Now suppose the result holds 
for some $n$. Then $h(n^+)=h(n)\cup\{C(A-h(n))\}=\{a_0,...,a_n,C(A-h(n))\}$. But this means $a_{n^+}=C(A-h(n))$, meaning $h(n^+)=\{a_0,...,a_n,a_{n^+}\}$, therefore completing 
the induction. Then $a_{n^+}$ is the sole element of $h(n^+)-h(n)=(h(n)\cup\{C(A-\{a_0,...,a_n\})\})-h(n)=\{C(A-\{a_0,...,a_n)\}$, that is, $a_{n^+}$ is the $C$-chosen 
element of $A-\{a_0,...,a_n\}$. Thus $a_0$ is the $C$-chosen element of $A$ and $a_{n^+}$ is the $C$-chosen element of $A-\{a_0,...,a_n\}$, meaning each $a_i$ is unique, 
and thus the set of all $a_i$ is a denumerable subset of $A$. 

Note that the second proof of Theorem 4.2 provided in the book can be expanded in essentially the same was as this exercise.

\section*{Section 6}

\subsection*{Theorem 6.1}
Here is why $\Sigma$ is closed under chain unions. For consider some nest $\mathcal{N}$ of elements of $\Sigma$. 
Then $\cup\mathcal{N}$ is a 1-1 function (by Proposition 1.4 of Chapter 6). 
Consider the domain $D$ of $\cup\mathcal{N}$. We show that the range $R$ of $\cup\mathcal{N}$ is $D\cup D'$ (which would mean that 
$\cup\mathcal{N}$ is a 1-1 function from some subset $S$ of $A$ onto $S\cup S'$, meaning $\cup\mathcal{N}\in\Sigma$). If 
$x\in R$, then, for some function $f\in\mathcal{N}$, $x\in\textrm{Ran}(f)$. 
But $\textrm{Ran}(f)=\textrm{Dom}(f)\cup\textrm{Dom}(f)'$. Also $\textrm{Dom}(f)\subseteq D$ and 
$\textrm{Dom}(f)'\subseteq D'$. Thus $\textrm{Dom}(f)\cup\textrm{Dom}(f)'\subseteq D\cup D'$, meaning 
$x\in D\cup D'$.
On the other hand, if $x\in D\cup D'$, either $x\in D$ or $x\in D'$. If 
$x\in D$, then for some function $f\in\Sigma$, $x\in\textrm{Dom}(f)$. But $f$ is 1-1 onto $\textrm{Dom}(f)\cup\textrm{Dom}(f)'$, meaning 
$x\in\textrm{Ran}(f)$, and thus $x\in R$. If $x\in D'$, then for some $f\in\mathcal{N}$ and some $y\in\textrm{Dom}(f)$, $x=\pi(y)$. But 
$f$ is 1-1 onto $\textrm{Dom}(f)\cup\textrm{Dom}(f)'$, meaning $x\in\textrm{Ran}(f)$, and thus $x\in R$. Hence $R=D\cup D'$, 
meaning $\cup\mathcal{N}\in\Sigma$.

It is remarked that $(M\cup M')\times (M\cup M')=(M\times M)\cup (M\times M')\cup (M'\times M)\cup (M'\times M')$. This is by 
Exercise 7.1 of Chapter 2.

It is remarked that $\phi\cup\psi$ is 1-1 from $M\cup M'$ onto $(M\times M)$. Here is why this is. We already know 
$\phi$ is a 1-1 map from $M$ onto $M\times M$, and $\psi$ is a 1-1 map from $M'$ onto $N$. Because 
$N=\cup\{M\times M',M'\times M,M'\times M'\}$, $N$ and $M\times M$ are disjoint (because $M$ and $M'$ are disjoint). This is 
why $\phi\cup\psi$ is 1-1.

\subsection*{Theorem 6.4}
Here is why $\Sigma$ is closed under chain unions. For consider some chain $C$ of elements of $\Sigma$. Then $\cup C$ is a function, 
in fact a 1-1 map (by Proposition 1.4 of Chapter 6). Let $D$ be the domain of $\cup C$. 
We need to show that the range of $\cup C$ is $D\times D$ (which would mean that 
$\cup C\in\Sigma$). On the one hand, suppose that $x\in\textrm{Ran}(C)$. Then, for some $f\in C$, $x\in\textrm{Ran}(f)$. 
But $\textrm{Ran}(f)=\textrm{Dom}(f)\times\textrm{Dom}(f)$, meaning $x\in D\times D$. On the other hand, suppose that 
$x\in D\times D$. Then $x=\langle y,z\rangle$ for some $y,z\in D$. Then, for some functions 
$f$ and $g$ (both in $C$), $y\in\textrm{Dom}(f)$ and $z\in\textrm{Dom}(g)$. But, $f$ and $g$ are comparable (because $C$ is a chain).
Without loss of generality, suppose that $f\subseteq g$. Then $y$ and $z$ are both in the domain of 
$g$, and because $g$ is 1-1 onto $\textrm{Dom}(g)\times\textrm{Dom}(g)$, $x\in\textrm{Ran}(g)$, and thus 
$x\in R$. Hence $D\times D=R$, meaning $\cup C\in\Sigma$.

Note why, if we assumed that $A-M$ were not smaller than $M$, we could conclude that $A-M$ would include a subset $M'$ of the same 
size as $M$. This is because we are assuming the axiom of choice, and thus the trichotomy of cardinals holds, meaning either 
$A-M\preceq M$ or $M\preceq A-M$. Because $A-M\prec M$ doesn't hold, this means that $M\preceq A-M$ holds. 

It should be noted that $M$ must be infinite (or else we could not use Theorem 6.3). For $M$ is the domain of some function 
$\phi$ that is 1-1 from $M$ onto $M\times M$, which means $M$ must be infinite.

It is concluded from $M\cong A$ (say $\phi$ places $M$ into a 1-1 correspondence with $A$)
and $M\cong M\times M$, that $A\cong A\times A$. Here is why this is. 
First, $M\times M\cong M\times A$. This is because the set of all ordered pairs of the form 
$\langle \langle x,y\rangle,\langle x,\phi(y)\rangle\rangle$ where $x\in M$ and $y\in M$ is a 1-1 correspondence between 
$M\times M$ and $M\times A$. In a similar fashion, $M\times A\cong A\times A$. Thus $M\times M\cong A\times A$.

\section*{Section 7}

\subsection*{Notes}
When we are talking about isomoprhims from an ordinal $\alpha$ onto a set $S$ of ordinals (such as in Lemma 7.2 and Proposition 7.3), it 
seems that the well ordering of $S$ must be the well ordering of $On$.

In part 2 of Proposition 7.4, it is concluded that $\alpha <d$. Here is why this is: We know that $S\cong\alpha,c\cong S$, and $c\prec d$. Also, 
by Proposition 7.3, $\alpha\le d$. If $\alpha =d$, then $c\cong d$, contrary to the fact that $c\prec d$. Thus $\alpha<d$. 

In the same proof as the preceeding paragraph, it was concluded that $c\le d$. Here is why this is: Suppose 
$c>d$. Then, as we have already shown in part (2) of the proof, $c\succ d$, contrary to the assumption that 
$c\prec d$.

There is in fact an alternative method of showing $c<d$ (in the second paragraph of Proposition 7.4 (2)): Suppose that 
$c\ge d$. Recall that $\alpha<d$ (meaning $\alpha<c$) and $c\cong\alpha$. 
But $c$ is a cardinal, meaning it is not the same size as any lesser ordinal, and 
we have $c\cong\alpha$, an ordinal less than $c$. Thus this is a contradiction, meaning it must be that $c<d$.

\subsection*{Exercise 7.1}
(a) Suppose $L$ is a lower section of a well ordered system $\langle A,\le\rangle$ and $\phi$ is an isomoprhism (with respect to $\le$) from 
$L$ into $A$. Say that an element $x$ of $L$ moves backwards if $x>\phi(x)$. Suppose some $x\in L$ moved backwards. Because $L$ is a lower section of 
$A$, this would mean that $\phi(x)\in L$.
Then, because $\phi$ is an isomorphism, $\phi(x)>\phi(\phi(x))$. Thus $\phi(x)$ also moves backwards. But this means that there is no least element of 
$L$ which moves backwards, contrary to the fact that $L$ is well ordered (because it is a subclass of $A$, which is well ordered). Hence no element of 
$x$ moves backwards, meaning for every $x\in L$, $x\le\phi(x)$.

(b) Here is how the above result yields both Theorem 2.1 of Chapter 6, and Lemma 7.2 of this chapter. In the case of Lemma 7.2, $\alpha$ is a lower section of 
$On$ (because every ordinal is a lower section of $On$), and $On$ is well ordered, and $\phi$ is an isomorphism (with respect to $\le$) from $L$ into $On$, meaning, 
by the above result, that for all $x\in\alpha$, $x\le\phi(x)$. Now, recall that Theorem 2.1 of Chapter 6 stated that, for an isomorphism $\phi$ from a 
well ordered system $\langle A,\le\rangle$ into itself, no element $a$ moves backwards under $\phi$. This follows from part (a) of this exercise, as 
$\phi$ is assumed to be an isomorphism from the whole of $A$, which itself is a lower section of $A$.

\section*{Section 8}

\subsection*{Notes}
Here is why we can do cardinal addition, multiplicaiton, and exponentiation without the axiom of choice. 
We know that $On$ can be properly well ordered, and so can $On^2$. This explains why, for any cardinals 
$c_1$ and $c_2$, $c_1+c_2$ and $c_1\cdot c_2$ will have cardinality, even without assuming the axiom of choice.

Now, the above paragraph explains why every function from an ordinal to another ordinal can be well ordered, as such 
a function is a subset of $On^2$ (this means the function is properly well orderd by Proposition 1.2 of Chapter 4). 
Let $S$ be the set of all functions from $d$ to $c$ for any cardinals $c$ and $d$. We can define a proper well ordering $\le$ on 
$S$ in the following way: for any functions $f$ and $g$ in $S$, 
$f\le g$ iff $f\subseteq g$ or, neither $f\subseteq g$ nor $g\subseteq f$ and 
the least ordered pair $\langle x,y\rangle$ in $f$ but not in $g$ is less than the least ordered pair 
$\langle p,q\rangle$ in $g$ but not in $f$. This can be seen to be a well ordering. Hence there really is some 
cardinal $c^d$. 

\subsection*{Theorem 8.6}
It is remarked that, if $\alpha$ and $\beta$ are both less than $c$ (where $c$ is an infinite cardinal), then $P(\alpha,\beta)<c$. Here is why 
this is.

\begin{lemma}
Every finite ordinal $\alpha$ is equal to $P(\beta,\gamma)$ for some finite ordinals $\beta,\gamma$.
\end{lemma}
\begin{proof}
By induction. $0=P(0,0)$ obviously holds, meaning the base case holds. Suppose, for all ordinals $\epsilon$
less than some ordinal $\alpha$, the $P(\beta,\gamma)$ for which $\epsilon=P(\beta,\gamma)$ is such that 
$\beta$ and $\gamma$ are both finite. Suppose that the $P(\beta,\gamma)$ for which 
$\alpha=P(\beta,\gamma)$ was such that at least one of $\beta,\gamma$ is infinte.
Now, $\alpha$ must be the successor of some ordinal; 
say it is the successor of $P(\delta,\epsilon)$. By assumption, $\delta$ and $\epsilon$ are both finite. 
Without loss of generality, suppose that $\textrm{max}(\delta,\epsilon)=\delta$. Then 
$P(\delta,\epsilon)<P(\delta+1,\epsilon)$. But also 
$P(\delta+1,\epsilon)<P(\beta,\gamma)$. But this contradicts the fact that 
$P(\beta,\gamma)$ is the successor of $P(\delta,\epsilon)$. Thus neither of $\beta$ or $\gamma$ can 
be infinite, meaning both $\beta$ and $\gamma$ must be finite.
\end{proof}

\begin{lemma}
For all finite ordinals $\alpha,\beta$, $P(\alpha,\beta)$ is a finite ordinal.
\end{lemma}
\begin{proof}
This is proven noting the following fact about a relation $R(x,y)$ on natural 
numbers $x$ and $y$. If:\newline
$C_1$ $R(0,0)$ holds;\newline
$C_2$ if $R(x,y)$ holds, then both $R(x+1,y)$ and $R(x,y+1)$ hold;\newline
then $R(x,y)$ holds for all natural numbers $x$ and $y$.

$C_1$ $P(0,0)$ is finite, as $P(0,0)=0$.

$C_2$ Suppose $P(x,y)$ is a natural number for finite ordinals $x$ and $y$. 
There are three cases: $\textrm{max}(x,y)=y$ and $x\ne y$; $x=y$;
$\textrm{max}(x,y)=x$ and $x\ne y$. We give the proof of the first case - the other three 
cases are similar.

So suppose $\textrm{max}(x,y)=y$ and $x\ne y$. If $x=y-1$, then 
the successor of $\langle x,y\rangle$ is $\langle x+1,0\rangle$, meaning 
$P(x+1,0)$ is finite, and an obvious induction argument 
shows that $\langle x+1,z\rangle $ for $z\le x+1$ is such that 
$P(x+1,z)$ is finite. Hence $P(x+1,y)$ is finite (because $y=x+1$). 
Then $\langle 0,y+1\rangle$ is the successor of $\langle x+1,y\rangle$, meaning 
$P(0,y+1)$ is finite. An obvious induction argument then shows that $P(z,y+1)$ for all 
$z\le y$ is finite, therefore meaning $P(x,y+1)$ is also finite. Hence both 
$P(x+1,y)$ and $P(x,y+1)$ are finite.

If $x<y-1$, then $\langle x+1,y\rangle$ is the successor of $\langle x,y\rangle$, meaning 
$P(x+1,y)$ is finite. An obvious induction argument shows that $P(z,y)$ is finite for all 
$z<y$, meaning $P(y-1,y)$ is finite. Then a similar argument as in the preceeding paragraph shows that 
$P(x,y+1)$ is finite.
\end{proof}

Then, going back to Theorem 8.6, $P(\alpha,\beta)$ must be finite if $\alpha$ and $\beta$ are both finite, meaning 
$P(\alpha,\beta)<c$.

\subsection*{Exercise 8.1}
Suppose that $a$ and $a'$ are the same size, and that $b$ and $b'$ are the same size, and $a$ and $b$ are disjoint. Suppose 
$f$ is a 1-1 function from $a$ onto $a'$ and $g$ is a 1-1 function from $b$ onto $b'$.
Then the function $\phi$ from $a\cup b$ onto $a'\cup b'$ such that $\phi(x)=f(x)$ if 
$x\in a$ and $\phi(x)=g(x)$ if $x\in b$ is 1-1, meaning $a\cup b$ and $a'\cup b'$ are of the same size.

\subsection*{Exercise 8.2}
Suppose that $a$ and $a'$ are the same size, and $b$ and $b'$ are the same size. Suppose $f$ is a 1-1 function 
from $a$ onto $a'$ and $g$ is a 1-1 function from $b$ onto $b'$.
Then the function $\phi$ from $a\times b$ onto $a\times b'$ such that 
$\phi(\langle x,y\rangle)=\langle x,g(y)\rangle$ where $x\in a$ and $y\in b$ is 1-1. Then the function 
$\psi$ from $a\times b'$ onto $a'\times b'$ such that 
$\psi(\langle x,y\rangle)=\langle f(x),y\rangle$ where $x\in a$ and $y\in b'$ is 1-1. Thus 
$a\times b$ is the same size as $a'\times b'$. 

\subsection*{Exercise 8.3}
Suppose $a,b,$ and $c$ are cardinals.

(a) We show that $a\times b$ is the same size as $b\times a$, which would mean that 
$a\cdot b=b\cdot a$. Let $\phi$ be the function from 
$a\times b$ onto $b\times a$ such that $\phi(\langle x,y\rangle)=\langle y,x\rangle$ where 
$x\in a$ and $y\in b$. It is easy to show that $\phi$ is 1-1.

(b) We show that $a\times (b\times c)$ is the same size as $(a\times b)\times c$, which would 
mean that $a\cdot (b\cdot c)=(a\cdot b)\cdot c$. Let $\phi$ be the function from 
$a\times (b\times c)$ onto $(a\times b)\times c$ such that 
$\phi(\langle x,\langle y,z\rangle\rangle)=\langle\langle x,y\rangle,z\rangle$ where 
$x\in a$, $y\in b$, and $z\in c$. It is easy to show that $\phi$ is 1-1.

(c) $a\cdot 0=0$ as the cardinality of $a\times 0$ is 0.

(d) $a\cdot 1=a$ as $a\times 1$ is just the set of all ordered pairs 
$\langle x,0\rangle$, where $x\in a$, and that can be put into a 1-1 correspondence with 
$a$.

(e) If $a\le b$, then $a\times c$ can be placed into a 1-1 correspondence with the set of all ordered pairs 
$\langle x,y\rangle$ where $x\in b$ and $x\in a$, and $y\in c$, and this is a subset of $b\times c$, meaning 
$a\cdot c\le b\cdot c$. 

$\le$ cannot be replaced by $<$ here. For instance, by Corollary 8.7, for any infinite cardinal $c$, 
if $a<b\le c$, then $a\times c=c=b\times c$.

\subsection*{Exercise 8.4}
The definition of cardinal exponentiation reads: \textquote{For cardinals $c$ and $d$, by $c^d$ is meant the cardinality of 
the set of all functions from $d$ to $c$.} I was unclear if \textquote{to} meant \textquote{into} or \textquote{onto} here. 
I am assuming it meant \textquote{into}. My reasoning is that this makes exponentiation work out the same was as normal exponentiation. For 
instnace, when interpreting exponentiation with \textquote{into}, $2^2=4$, as the set of all functions from $2$ into $2$ just contains 
the functions $\{\langle 0,1\rangle,\langle 1,0\rangle\},\{\langle 0,0\rangle,\langle 1,1\rangle\}$,
$\{\langle 1,1\rangle,\langle 0,1\rangle\},$ and $\{\langle 0,0\rangle,\langle 1,0\rangle\}$, that is, it contains 
four functions.

Suppose $a,b,$ and $c$ are cardinals. Let $A$ be the set of all functions from 
$a$ to $b$, and $S$ be the set of all functions from $A$ to $c$. Then $(a^b)^c$ is the cardinality of $S$. Also, 
$a^{(b\times c)}$ is the cardinality of the set of all functions from $a$ to $b\times c$. For any function 
$f\in S$, let $\phi(f)=g$, where $g$ is the set of all ordered pairs
$\langle x,\langle y,z\rangle\rangle$ where $\langle x,y\rangle\in\textrm{Dom}(f)$ and 
$f(\langle x,y\rangle)=z$. It can be seen that $\phi$ is 1-1 from $S$ onto  
the set of all functions from $a$ to $b\times c$, and thus $(a^b)^c=a^{(b\times c)}$
(they have the same cardinality). 

\subsection*{Exercise 8.5}
Suppose $c$ is an infinite cardinal and $d$ is a non-zero cardinal $\le c$. This means 
$d\cdot c\le c\times c$, as $d\times c\subseteq c\times c$. Also, by Theorem 8.6, 
$c\cdot c=c$, meaning $d\cdot c\le c$. But, because $d$ is non-zero, 
$c\preceq d\times c$. Hence $d\cdot c\cong c$. But, because $d\cdot c$ and $c$ are both cardinals, this 
means $d\cdot c=c$.

\end{document}
